\section*{结论}\addcontentsline{toc}{section}{结论}
本次公开修订将研究重点从复刻旧数值扩展到可靠估计和可解释比较。保留历史证据的同时，采用解析梯度与确定性多起点求解DEC，在上证50特征构造前排除已确认的休市行，并新增不含EPU的XGB-AR对照。沪深300中HAR的MSE、MAE最低，GARCH的QLIKE最低；DEC相对XGB-Market降低QLIKE和MAE，但MSE略高。

新增消融表明，三市场加入MIDAS-EPU后相对XGB-AR均降低QLIKE、提高MSE，中国两市场的MAE下降而纳斯达克上升。由此应将结论表述为当前设定下依赖损失函数的增量预测表现，不是EPU普遍改善预测或产生结构性因果效应。HAC和Holm调整提供探索性支持，不能消除回顾性分析、非线性估计及规格选择的推断限制。

纳斯达克验证最优规格从旧版12个月Beta改为12个月指数Almon。旧Beta局部解的训练目标更高，却偶然具有更低的验证误差；修订求解后按相同验证规则重选规格。相对旧DEC，修订测试MSE下降而QLIKE上升。这说明方法修正的价值在于估计规则清楚、数值更稳定和证据可追溯，并不要求所有测试指标都改善。

共同1859个日期上，主设置的DEC区间得分为10.1995，高于HAR的9.6973、XGB-Market的9.8534和XGB-AR的9.8793。DEC没有改善这些基准的区间质量。主设置覆盖诊断不拒绝不构成有效性证明；95\%、窗口120的DEC仍有违约独立性及条件覆盖拒绝结果。区间只针对观测代理及统一残差构造。

数值追溯已经精确重现6888条历史XGBoost预测；旧DEC历史环境仍不能唯一还原，新DEC在两套数值依赖中的最大预测差异小于$10^{-7}$。多起点不保证全局最优，边界权重的经济解释应克制。描述统计已使用3941条真实有效收益率；其他日期质量、EPU实时发布与修订、对数反变换以及独立样本表现仍待核验。未来应以这些问题及确认性留出评估为重点，不将当前结果包装成已验证的实时预测系统。
